\documentclass[../main.tex]{subfiles}

\begin{document}

\section*{第1問}

{\bf [解]}
$Q(x)=0$は重解を持たないと仮定する．$Q(x)=A(x-\alpha)(x-\beta)$（$\alpha\ne\beta$，$A\ne0$）とおく．

$P(x)$は$Q(x)$で割り切れないから，$P(x)$は$(x-\alpha)$，$(x-\beta)$の少なくとも一方を因数に持たない．

いずれの場合も，$P(x)$がその因数（例えば$(x-\alpha)$）を持たなければ，$P(x)^2$も$(x-\alpha)$を因数に持たず，$Q(x)$（$(x-\alpha)$を因数に持つ）で割り切れない．

これは，$P(x)^2$が$Q(x)$で割り切れるという仮定に矛盾する．よって$Q(x)$は重解を持つ．$\blacksquare$

\newpage

\section*{第2問}

{\bf [解]}
線分$AB$上の点は
\begin{align*}
\begin{pmatrix}x\\y\\z\end{pmatrix}=s\begin{pmatrix}2\\3\\0\end{pmatrix}+(1-s)\begin{pmatrix}0\\1\\2\end{pmatrix}\quad(0\le s\le1)
\end{align*}
とおける．この時，$k$は$(0\le k\le1)$に対し，
\begin{align*}
s\begin{pmatrix}2\\3\\0\end{pmatrix}+(1-s)\begin{pmatrix}0\\1\\2\end{pmatrix}=k\begin{pmatrix}5+t\\9+2t\\5+3t\end{pmatrix}
\end{align*}
なる$k$があることを示せば良い．
\begin{align*}
2s&=k(5+t)\\
2s+1&=k(9+2t)\\
-2s+2&=k(5+3t)
\end{align*}
\begin{align*}
\therefore\ t=-1，\quad k=\frac13，\quad s=\frac23
\end{align*}

だから，$t=-1$で交点を持ち，その座標は
\begin{align*}
\left(\frac43,\frac73,\frac23\right)．
\end{align*}

\newpage

\section*{第3問}

{\bf [解]}
$x\le0$の時，$f(x)=-\left(x+\dfrac12\right)^2+\dfrac14$，$x\ge0$の時，$f(x)=\left(x-\dfrac12\right)^2-\dfrac14$である．$x\le0$の時，$f'(x)=-2\left(x+\dfrac12\right)$だから，$x=-1$の接線は
\begin{align*}
y=x+1
\end{align*}
であるから，$x\ge0$での交点は$x=1+\sqrt2$である．

\begin{figure}[htb]
\centering
\begin{tikzpicture}[scale=1.4, >=stealth]
  \draw[->] (-1.8,-0.3) -- (3.2,-0.3+0.5) node[right] {$x$};
  \draw[->] (0,-0.6) -- (0,3.2) node[above] {$y$};
  \draw[domain=-1.5:0, smooth, thick] plot (\x,{-(\x+0.5)*(\x+0.5)+0.25});
  \draw[domain=0:2.42, smooth, thick] plot (\x,{(\x-0.5)*(\x-0.5)-0.25});
  \draw[domain=-1.3:2.6, smooth] plot (\x,{\x+1});
  \node at (-1,-0.6) {$-1$};
  \node at (2.42,-0.6) {$1+\sqrt2$};
\end{tikzpicture}
\caption{$y=f(x)$と$x=-1$における接線$y=x+1$で囲まれる部分（斜線部）}
\end{figure}

あって，求める面積$S$として
\begin{align*}
S&=\int_{-1}^0\left\{(x+1)+(x^2+x)\right\}dx+\int_0^{1+\sqrt2}\left\{(x+1)-(x^2-x)\right\}dx\\
&=\int_{-1}^0(x+1)^2\,dx+\int_0^{1+\sqrt2}(-x^2+2x+1)\,dx\\
&=\frac13+\Bigl[-\frac13x^3+x^2+x\Bigr]_0^{1+\sqrt2}\\
&=\frac13+(1+\sqrt2)+(1+\sqrt2)^2-\frac13(1+\sqrt2)^3\\
&=\frac{6+4\sqrt2}3．
\end{align*}

\newpage

\section*{第4問}

{\bf [解]}
$n$が素数で$n\ne3$の時，$n\not\equiv0\pmod3$であり，$n^2+2\equiv0\pmod3$となり，$n^2+2$（$>3$）は素数ではない．

$n=3$の時，$n^2+2=11$は素数．$n=2$の時，$n^2+2=6$は素数でない．したがって，示された．$\blacksquare$

\newpage

\section*{第5問}

{\bf [解]}
点$X$に対し，$\overrightarrow{AX}=\vec x$とおく．$s,t,u$を$0<s,t,u<1$を満たす実数として
\begin{align*}
\vec p=s\vec b，\qquad\vec q=t\vec c+(1-t)\vec b，\qquad\vec r=u\vec c
\end{align*}
とおける（$\vec b=\overrightarrow{AB}$，$\vec c=\overrightarrow{AC}$）．$\triangle PQR$の重心を$G$とする．
\begin{align*}
3\vec g&=(t+u)\vec c+(1+s-t)\vec b
\end{align*}

ここで，$\lambda=1+s-t$，$\mu=t+u$とおくと，$s,t,u\in(0,1)$が独立に動くことから，$(\lambda,\mu)$は頂点
\begin{align*}
(1,0),(2,0),(2,1),(1,2),(0,2),(0,1)
\end{align*}
を持つ六角形の内部（境界含まず）を動く．$\vec g=\dfrac\lambda3\vec b+\dfrac\mu3\vec c$だから，これは，各辺を$1:2$，$2:1$に内分する6点を頂点とする六角形の内部に対応する．

だから，下図斜線部（境界含まず）．

\begin{figure}[htb]
\centering
\begin{tikzpicture}[scale=2, >=stealth]
  \coordinate (A) at (1,1.7);
  \coordinate (B) at (0,0);
  \coordinate (C) at (1.6,0);
  \draw (A) -- (B) -- (C) -- cycle;
  \coordinate (P1) at ($(A)!1/3!(B)$);
  \coordinate (P2) at ($(A)!2/3!(B)$);
  \coordinate (Q1) at ($(B)!1/3!(C)$);
  \coordinate (Q2) at ($(B)!2/3!(C)$);
  \coordinate (R1) at ($(C)!1/3!(A)$);
  \coordinate (R2) at ($(C)!2/3!(A)$);
  \draw[dashed, fill=gray!25] (P1) -- (P2) -- (Q1) -- (Q2) -- (R1) -- (R2) -- cycle;
  \node[above] at (A) {$A$};
  \node[below left] at (B) {$B$};
  \node[below right] at (C) {$C$};
\end{tikzpicture}
\caption{重心$G$が動く範囲（斜線部，境界含まず）}
\end{figure}

\newpage

\section*{第6問}

{\bf [解]}
$F'(\theta)=\theta\cos(\theta+\alpha)$だから下表を得る．

\begin{center}
\begin{tabular}{c|ccccc}
$\theta$ & $0$ & & $\dfrac\pi2-\alpha$ & & $\pi/2$ \\
\hline
$F'$ & & $+$ & $0$ & $-$ & \\
\hline
$F$ & & $\nearrow$ & & $\searrow$ & \\
\end{tabular}
\end{center}

したがって，$F(\theta)$は$\theta=\dfrac\pi2-\alpha\quad\cdots①$で最大．
\begin{align*}
F(\theta)&=\Bigl[x\sin(x+\alpha)+\cos(x+\alpha)\Bigr]_0^{\pi/2-\alpha}\quad(\because①)\\
&=\left(\frac\pi2-\alpha\right)-\cos\alpha．
\end{align*}

\newpage

\end{document}
